\documentclass[11pt]{article}
\usepackage[preprint]{acl}
\usepackage{booktabs}
\usepackage{graphicx}
\usepackage{amsmath}
\usepackage{microtype}
\hypersetup{
  pdftitle={Language-Model Distributions as Tools for Computational Linguistics},
  pdfauthor={Akufor Agho Fru}
}

\title{Language-Model Distributions as Tools for Computational Linguistics}
\author{Akufor Agho Fru}

\begin{document}
\maketitle

\begin{abstract}
Computational linguistics uses probability distributions to rank documents,
identify languages, and describe lexical behavior. I test when
language-model distributions help these tools, against matched counts and
against the standard baselines for each task. The answer is mostly no. In
query-likelihood retrieval, a hybrid of DistilGPT-2 word probabilities and
counts beats a background estimated from only 1,000 sampled tokens (SciFact
nDCG@10 .650 versus .599). That setting is artificial, because query
likelihood already reads the full collection. The gain disappears at larger
budgets, and the standard baselines score higher on SciFact: collection
Dirichlet smoothing reaches .659 and BM25 .663, although a paired test of BM25
against the hybrid is not significant. LM-derived IDF and LM-based stop words do not beat
BM25. For Bosnian/Croatian and Indonesian/Malay identification, a causal model
trained from scratch reaches 77.94\% raw accuracy, against 76.24\% for a
unigram that shares its subword tokenizer and 82.19\% for GlotLID-M. On
Indonesian/Malay the raw unigram is better (92.59\% versus 89.67\%). After a
post-hoc development calibration, the causal model beats the unigram by 2.03
points and a post-hoc character 1-5 gram naive Bayes control by 2.47 points,
but GlotLID-M stays best (86.34\%). Stylometry and distinctive-term classification show no
reliable gain. The gains in next-word prediction and EIFFEL idiom completion align with
the language-modeling objective.
They show cloze ability, not linguistic utility beyond it. The dative result
failed its registered standalone test and passed only a post-hoc additive
test, consistent with Hawkins et al.\ (2020). Language-model
distributions are candidate measurements to test against task baselines, not
replacements for them.
\end{abstract}

\section{Introduction}

Many computational-linguistics tools need a distribution over words: a
background for retrieval smoothing, a reference rate for stylometry, a prior
for spelling correction, or a likelihood for language identification. Large
language models now give cheap estimates of such distributions. A natural
question is whether these estimates make the tools better than the counts and
baselines they would replace.

I answer this question with saved evaluations across retrieval, language
identification, stylometry, lexical distinctiveness, next-word prediction,
spelling correction, stop-word selection, existing human-judgment datasets,
and lexical distributional laws. The studies use
different models, corpora, and designs. I report each one as a separate
evidence block and do not pool their metrics. Three research questions
organize the evidence:

\begin{description}
\item[RQ1] Does a frozen language-model word distribution improve retrieval,
stylometry, and lexical tasks over matched empirical counts and over the
standard baselines for each task?
\item[RQ2] Does a contextual model improve close-language identification over
a matched subword unigram, a character n-gram model, and released
identifiers?
\item[RQ3] Do language-model probabilities add value on existing human-judgment
data, and does any gain go beyond the language-modeling objective?
\end{description}

The answer is conditional and mostly negative. An interpolated prior helps
retrieval only when the empirical background is very sparse, and the standard
retrieval baselines remain as good or better. In the planned raw evaluation, a
contextual classifier is mixed against a matched unigram and trails GlotLID-M.
After a post-hoc calibration, it beats the unigram and a post-hoc character
n-gram control, but still trails GlotLID-M. The
positive human-judgment results align with the cloze-style training
objective or repeat known findings.

The contributions are:
\begin{enumerate}
\item A matched comparison of frozen language-model word priors with sampled
counts at several data budgets, with the standard retrieval baselines reported
next to them.
\item A close-language identification comparison with raw results and a
post-hoc calibrated analysis, which adds a character n-gram control, so the
contextual gain is not measured only against a baseline that shares the
model's tokenizer.
\item Evidence that confirms the framing of \citet{trott2026distributional}
across several tasks. Gains appear where the task matches the training
objective, such as next-word prediction and idiom completion, and not where
the task needs more than the objective. This separation is not new.
\end{enumerate}
This is an evidence synthesis, not a new estimator. I make no claim that
language models replace empirical distributions.

\section{Related Work}

\paragraph{Smoothing and retrieval.} Query-likelihood retrieval ranks a
document by the probability of the query under a document language model
\citep{pontecroft1998}, and translation models extend it with word-to-word
probabilities \citep{bergerlafferty1999}. Unigram smoothing is central to this
approach. \citet{zhai2004smoothing} showed that the background
distribution and its parameters affect query-likelihood performance, so
interpolating counts with a background is a standard construction. BM25 is the
standard lexical baseline \citep{robertson2009bm25}. Neural language models
already improve lexical retrieval in stronger ways than a frozen background
prior. DeepCT learns contextual term weights \citep{daicallan2020deepct},
SPLADE learns sparse lexical expansions \citep{formal2021splade}, and
doc2query expands documents with predicted queries
\citep{nogueira2019doc2query}. \citet{dossantos2020beyond} rank passages by the
likelihood a GPT-2 model gives to the query, and \citet{zhuang2023} use
open-source LLMs as zero-shot query-likelihood models. My retrieval test is narrower.
It replaces only the background distribution of a Dirichlet query-likelihood
model with frozen word probabilities and asks whether that helps at small
reference budgets. I do not compare with these neural rankers.

\paragraph{Stylometry and lexical statistics.} Word-based authorship methods
compare observed frequencies with corpus statistics, as in Burrows's Delta
\citep{hoover2004testing}. \citet{eder2022boosting} showed that changing the
normalization denominator can improve attribution, and
\citet{peng2003authorship} attributed authors with character-level n-gram
language models. Frequency adjustment and
prior-based comparison are therefore not new. \citet{barlas2020crossdomain} used pretrained language models for
cross-domain attribution, and \citet{huang2025perplexity}
fine-tune one authorial language model per candidate. My stylometry study
instead uses one frozen global model as an external word prior. Informative
Dirichlet log odds is an established way to rank distinctive words
\citep{monroe2008fightin}. PMI and the t-score are standard association
baselines \citep{bouma2010collocation}. Human-judged collocation resources
exist \citep{kilgarriff2014collocation,collocaid2021}, but I could not use them
here. The first has no reuse terms I could verify, and the second has no
candidate-level negative judgments. I therefore use held-out next-word
prediction, which is a different task. Noisy-channel spelling correction
combines a candidate prior with an error channel \citep{brillmoore2000}.

\paragraph{Language models as baselines and estimators.} Language-model
probabilities have long served as measurements in psycholinguistics. Surprisal
theory links word probability to processing difficulty
\citep{hale2001earley,levy2008expectation}. Reading times grow with surprisal
\citep{smithlevy2013logarithmic,shain2024logarithmic}. Better language models
gave better reading-time fits up to a point
\citep{goodkind2018surprisal,wilcox2020predictive}, but larger transformers
fit worse \citep{ohschuler2023surprisal}. So better word prediction does not
always give a better measurement, which is also what my cross-entropy results
show. \citet{trott2026distributional} use LLMs as a confound control in
psycholinguistics. The model tests whether a human behavior could arise from
distributional exposure alone. In this framing, an LLM shows what
distributional statistics can achieve on a task. My third contribution fits
this framing: gains on objective-aligned tasks show what next-word statistics
can do, not a further linguistic ability.
\citet{zhong2025ratings} compare prompted LLM
ratings with human ratings and corpus measures. LM perplexity also tracks
learner proficiency. It falls as CEFR level rises in Swedish learner essays
\citep{munozsanchez2024perplexity}, and LLM perplexity features classify CEFR
levels across learner corpora \citep{stearns2026cefr}. My MERLIN pilot tests
the same idea on German learner texts. \citet{hawkins2020dais}
introduced the DAIS dative-alternation dataset and showed that neural
language-model probabilities predict human dative preferences. My dative pilot
repeats that test with newer models. Word probabilities from subword models
need a token-boundary correction, and I use the one of
\citet{pimentel2024word}. \citet{nguyenarehalli2026tokenizer} show that BPE
vocabulary size changes surprisal estimates, so the tokenizer is part of the
measurement.

\paragraph{Language identification.} Character n-gram models are the classical
approach to language identification \citep{cavnar1994ngram,dunning1994statistical},
as a survey of the field describes \citep{jauhiainen2019survey}.
Shared tasks on similar languages include Bosnian/Croatian/Serbian and
Indonesian/Malay \citep{zampieri2014dsl}. In the third DSL shared task, systems
built on character n-grams with SVMs, logistic regression, or n-gram language
models won, and neural entries were not competitive \citep{malmasi2016dsl}.
In that task, a character 7-gram language model reached 88.45\% accuracy and a
character convolutional LSTM 78.45\% \citep{cianflone2016ngram}.
Perplexity-based identification of similar languages is established
\citep{gamallo2017perplexity}. UniLID identifies languages from tokenizer
unigram distributions \citep{meister2026unilid}, so the matched subword
unigram in my study is a UniLID-style classifier. Released identifiers include
fastText lid.176 \citep{joulin2017bag,joulin2016fasttextzip}, GlotLID-M
\citep{kargaran2023glotlid}, and OpenLID \citep{burchell2023openlid}.
OpenLID-v3 targets closely related languages, including
Bosnian/Croatian/Serbian \citep{fedorova2026openlidv3}. I did not evaluate
OpenLID or OpenLID-v3.

\paragraph{Statistical laws.} Quantitative laws of word use predate language
models. \citet{takahashi2017neural,takahashi2019scaling} already measured
Zipf, Heaps, and long-range scaling properties in neural language models and
their generated text. \citet{mikhaylovskiy2025zipf} measured Zipf and Heaps
behavior in LLM-generated text across temperature settings. Statistical
accounts of Zipf remain contested, and an absolute fit needs explicit model
comparison and goodness-of-fit tests \citep{piantadosi2014zipf,clauset2009powerlaw}.
\citet{kobayashi2018taylor} studied Taylor scaling in linguistic sequences.
My law measurements therefore describe two small models and two domains. They
make no new claim about Zipf, Heaps, Taylor, or temperature effects.

\section{Data, Models, and Licenses}

Table~\ref{tab:data} lists every dataset and model used in the evidence
blocks, with its source and license. I used each resource for research only
and do not redistribute dataset text with this paper. EIFFEL and MERLIN are licensed CC BY-SA 4.0, and FLORES-200 is CC BY-SA
4.0. I attribute them here and in the run records. Any derived data released
from these sources will use the same license. The Blog Authorship Corpus is
used under its noncommercial research terms. BEIR redistributes SciFact,
NFCorpus, and FiQA but does not grant upstream rights, so the table gives the
terms of each original source. NFCorpus is free for academic use only, and
the FiQA question and answer data are for noncommercial use only.

\begin{table*}[t]
\centering
\small
\begin{tabular}{p{4.2cm}p{5.4cm}p{4.6cm}}
\toprule
Resource & Use in this paper & License \\
\midrule
SciFact \citep{wadden2020scifact}, via BEIR \citep{thakur2021beir}
 & Retrieval, LM-IDF, and stop-word studies
 & Claims CC BY 4.0, abstracts ODC-By 1.0 \\
NFCorpus \citep{boteva2016nfcorpus}, via BEIR
 & Retrieval, LM-IDF, and stop-word studies
 & Free for academic use only \\
FiQA-2018 \citep{maia2018fiqa}, via BEIR
 & LM-IDF and stop-word studies
 & QA data noncommercial use only \\
WikiText-103 and WikiText-2 \citep{merity2017pointer}
 & LM contexts, empirical background, lexical laws
 & CC BY-SA 3.0 and GFDL \\
Blog Authorship Corpus \citep{schler2006effects,blogcorpus}
 & Stylometry, distinctive terms
 & Noncommercial research use \\
Project Gutenberg novels \citep{gutenberg}
 & Stylometry, next-word, spelling, lexical laws
 & Public domain in the USA \\
DAIS \citep{hawkins2020dais}
 & Dative-preference pilot
 & CC BY 4.0 \\
EIFFEL \citep{eiffel,eiffeldata}
 & French idiom-completion pilot
 & CC BY-SA 4.0 \\
MERLIN \citep{merlin2014,merlindata}
 & German learner proficiency pilot
 & CC BY-SA 4.0 \\
SWORDS \citep{lee2021swords}
 & Lexical-substitution pilot
 & CC BY 3.0 US \\
French Lexical Network idiomaticity data \citep{liu2024idiom,fln}
 & Idiomaticity-class pilot
 & CC BY 4.0 \\
FLORES-200 \citep{nllb2022,goyal2022flores}
 & Language-identification dev and devtest
 & CC BY-SA 4.0 \\
Bosnian, Croatian, Indonesian, and Malay Wikipedia, 2023-11-01 \citep{wikimedia2023}
 & Language-identification training text
 & CC BY-SA 3.0 and GFDL \\
\midrule
DistilGPT-2 \citep{sanh2019distilbert,distilgpt2}
 & Frozen word priors, LM-IDF, lexical laws
 & Apache-2.0 \\
GPT-2 \citep{radford2019gpt2}
 & Lexical laws
 & MIT \\
Qwen3.5-2B-Base and Qwen3.5-4B-Base \citep{qwen35}
 & Human-judgment pilots
 & Apache-2.0 \\
fastText lid.176 \citep{joulin2017bag}
 & Language-identification comparator
 & CC BY-SA 3.0 \\
GlotLID-M \citep{kargaran2023glotlid}
 & Language-identification comparator
 & Apache-2.0 plus notices \\
\bottomrule
\end{tabular}
\caption{Datasets and models, their use, and their licenses as stated by the
official source. Model revisions and dataset checksums are in the run
records.}
\label{tab:data}
\end{table*}

\section{Statistical Framework}

\paragraph{Shared prior construction.} Let $V$ be the fixed evaluation
vocabulary, $c_w$ the count of word $w$ in the task's reference sample, and
$q_{\mathrm{LM}}$ the frozen model's word probability normalized over $V$.
I compare the empirical distribution $q_E(w)=c_w/\sum_{v\in V}c_v$, the
language-model distribution $q_{\mathrm{LM}}$, and the hybrid
\begin{equation}
 q_H(w)=\frac{c_w+\beta q_{\mathrm{LM}}(w)}
 {\sum_{v\in V}c_v+\beta}.
\end{equation}
When a reference sample has zero counts for vocabulary items, empirical
probabilities use a $0.5$ additive count before normalization. The hybrid's
mixing mass $\beta$ and downstream smoothing parameters are selected on
training or development data only.

I use the frozen DistilGPT-2 checkpoint \citep{sanh2019distilbert,distilgpt2}
at a pinned revision recorded in run metadata. For context set $C$,
$q_{\mathrm{LM}}(w)$ is the normalized mean of boundary-corrected whole-word
probabilities, $|C|^{-1}\sum_{c\in C}P_{\mathrm{LM}}(w\mid c)$. I average over
32 independent WikiText-103 test contexts \citep{merity2017pointer}. The
context sample is independent of all task references. The fixed vocabulary
contains 300 shared stylometry features, retrieval query words, the blog
reference vocabulary, or the 20,000 most frequent novel reference words,
depending on task. Vocabulary construction uses no test labels or test
relevance judgments. Pair probabilities are normalized over each fixed
candidate set. Word priors are normalized over their task vocabulary, and
spelling posteriors are normalized across each case's candidate list.

\paragraph{Stylometry.} I use 1,000-token samples from the Blog Authorship
Corpus \citep{schler2006effects} and public-domain novels from Project
Gutenberg. Query documents and known-author references are disjoint texts from
the same author set. For each corpus, I estimate $q_E$ from a random sample of
1,000, 10,000, or 100,000 reference tokens using seeds 17, 29, and 43, where
feasible. I also evaluate a single full-reference estimate. The reference
splits contain 58,391 eligible words for blogs and 73,922 for novels, so the
100,000-token condition is infeasible for both and is omitted.
Document features are smoothed log rate residuals,
\begin{equation}
 r_w(d)=\log\frac{c_w(d)+\alpha q(w)}{N_d+\alpha}-\log q(w),
\end{equation}
where $N_d$ is the document length. I use mean absolute distance from each
query to every reference. The primary metric is global pairwise AUC, with
top-one accuracy as a descriptive measure. Empirical, LM, and hybrid methods
share each vocabulary and distance rule. Smoothing $\alpha$ and hybrid $\beta$
are selected by leave-one-reference-out attribution among training authors.
Burrows's Delta is included as a full-data reference baseline.

\paragraph{Retrieval.} I evaluate SciFact and NFCorpus from the BEIR archives.
Document counts are drawn uniformly without replacement at the same three
budgets and seeds, or use the full corpus. For each query, I rank all
documents by the query log likelihood under a Dirichlet-smoothed document
model with $q$ as its background distribution:
\begin{align}
 P(w\mid d)&=\frac{c(w,d)+\mu q(w)}{|d|+\mu},\\
 s(q,d)&=\sum_{w\in V}c(w,q)\log P(w\mid d).
\end{align}
The Dirichlet mass $\mu$ and hybrid $\beta$ are tuned on the collection's
development judgments. Test metrics are nDCG@10 and MAP@100.

I also report the standard baselines. These are full-index Dirichlet query
likelihood with the collection background, Jelinek-Mercer collection
smoothing, and BM25, together with uniform and additive smoothing and an
empirical WikiText background mixed with collection frequencies. The
smoothing parameters of the Dirichlet, Jelinek-Mercer, uniform, additive, and
WikiText-mix methods, and the hybrid's $\mu$ and $\beta$, are selected on
development judgments. BM25 uses fixed $k_1=1.2$ and $b=0.75$ and was not
tuned. This difference favors the hybrid. All methods use the same official
test queries. The 1,000-token background setting is artificial. Query likelihood
already reads the full collection to build document models, so in practice
the background comes from the full collection for free. The sampled budgets
isolate how a prior behaves when counts are scarce, and the standard
baselines show what a practitioner would actually use. Their
separate family of 40 contrasts is labeled exploratory. A separate paired test compares saved BM25
rankings with the 1,000-token hybrid on the same queries. It uses 10,000
paired query bootstrap resamples and a 10,000-sample two-sided sign-flip test,
with Holm correction over the two collections. I checked the five
selected controls with \texttt{ir-measures} 0.4.3 \citep{macavaney2022irmeasures}
from saved rankings. All 12,460 per-query metric values agreed within
$2.22\times10^{-16}$ using linear-gain nDCG and the saved rank order for ties.

\paragraph{Lexical distinctiveness.} I compare the adult Student and
Technology job-label groups in the Blog Authorship Corpus. Reference,
development, and test authors are disjoint. After short-reference filtering,
the reference split contains 72 Student and 76 Technology authors, with 40
development and 40 test authors. For each reference-token budget and seed,
informative Dirichlet log odds with prior mass 1,000 ranks terms by absolute z
score. A logistic classifier trained on development authors predicts
test-author labels from the top 100 terms. Accuracy is primary, with macro-F1
descriptive.

\paragraph{Next-word prediction.} I use 2,000 held-out adjacent-word cases from
public-domain novels, clustered by 20 book authors. Each case contains the true
next word and four distractors from the same reference-frequency bin.
Candidate sets are fixed across methods. The LM conditions only on the
preceding word. I compare it with sampled bigram counts, their hybrid, a
smoothed bigram backoff, a uniform prior, a full-reference frequency control,
PMI, and the t-score. The backoff probability is proportional to the observed
pair count plus 0.5 times the reference unigram probability, divided by the
pair count total plus 0.5, then normalized over the fixed candidate set.
Top-one accuracy is primary, with mean reciprocal rank (MRR) descriptive. This
task is the LM training objective itself. It is not an evaluation against human
collocation labels.

\paragraph{Spelling correction.} I generate 2,000 single-edit typos per seed
from held-out novel text, for three seeds. Each candidate list contains the
intended word and up to 24 words in the fixed training lexicon that are one
edit from the observed typo. Every prior uses the same candidate list and
uniform single-edit error channel. I compare top-one accuracy and MRR. This
controlled check does not measure naturally occurring typos.

\paragraph{Inference.} Retrieval outcomes are first averaged across the three
reference-sampling seeds for each query. Paired query bootstrap intervals and
two-sided sign-flip tests then compare methods. To test data scarcity, I also
compare each LM or hybrid benefit at the smallest and largest feasible sample
budgets using the same queries or author clusters. Intervals for the new tasks
resample authors, averaging seed-level outcomes within each author. Their
paired sign-flip tests use 10,000 samples. Stylometry uncertainty uses 2,000
author-cluster bootstrap draws, resampling the three reference seeds and author
clusters, and 10,000 author-cluster sign-flip permutations. Holm correction
covers the original 50 planned retrieval and stylometry tests plus 12 new-task
tests. Six post-hoc bigram-control contrasts are also reported, with adjusted
results across all 68 tests. The separate retrieval-control family contains 40
exploratory contrasts. I report held-out cross-entropy for each task's own
probability target, not as a substitute for task performance. Token-level
cross-entropy is used for stylometry, retrieval, and distinctiveness,
candidate-conditional cross-entropy for next-word prediction, and
channel-combined candidate-posterior cross-entropy for spelling.

\section{Experimental Design}

The design changes only the reference prior and its tuned smoothing
parameters. The task examples, test judgments, query lengths, feature
vocabularies, and evaluation metrics remain fixed within each comparison.
Retrieval uses 300 SciFact test queries and 323 NFCorpus test queries. The
stylometry comparisons use the existing held-out author splits. All language
model word scores were computed once and reused. Downstream sweeps do not
rerun the language model.

For retrieval, each combination of collection, token budget, and method is
evaluated on identical test queries. For stylometry, paired query-reference
distances are computed from each method's feature representation, and
uncertainty resamples author clusters to respect shared authors across
queries and references. The per-query data, seed-level results, test
statistics, and environment metadata are included with each run.

\section{Results}

\subsection{RQ1: Retrieval}

\begin{table*}[t]
\centering
\small
\begin{tabular}{llrrrr}
\toprule
 & & \multicolumn{2}{c}{SciFact} & \multicolumn{2}{c}{NFCorpus} \\
Tokens & Prior & nDCG@10 & MAP@100 & nDCG@10 & MAP@100 \\
\midrule
1,000 & Empirical & .599 & .562 & .290 & .130 \\
      & LM         & .642 & .605 & .306 & .142 \\
      & Hybrid     & .650 & .612 & .308 & .140 \\
10,000 & Empirical & .642 & .601 & .299 & .136 \\
       & LM         & .642 & .605 & .306 & .142 \\
       & Hybrid     & .631 & .592 & .308 & .141 \\
100,000 & Empirical & .653 & .609 & .298 & .135 \\
        & LM         & .642 & .605 & .306 & .142 \\
        & Hybrid     & .637 & .594 & .299 & .136 \\
Full & Empirical & .656 & .615 & .300 & .136 \\
     & LM         & .642 & .605 & .306 & .142 \\
     & Hybrid     & .656 & .615 & .301 & .137 \\
\bottomrule
\end{tabular}
\caption{Query-likelihood retrieval with each background prior. Mean nDCG@10
and MAP@100. Sampled-reference results average three seeds. The full reference
is one fixed estimate. Table~\ref{tab:retrieval-controls} gives the standard
baselines.}
\label{tab:retrieval}
\end{table*}

\begin{table}[t]
\centering
\small
\begin{tabular}{lrr}
\toprule
Method (nDCG@10) & SciFact & NFCorpus \\
\midrule
Hybrid prior, 1,000 tokens & .650 & .308 \\
\midrule
BM25 & .663 & .305 \\
Collection Dirichlet & .659 & .298 \\
Jelinek-Mercer & .657 & .293 \\
WikiText background mix & .661 & .305 \\
Uniform Dirichlet & .543 & .279 \\
Additive uniform & .495 & .278 \\
\bottomrule
\end{tabular}
\caption{Standard retrieval baselines next to the best LM-based result. BM25,
collection Dirichlet, and Jelinek-Mercer use the full BEIR index. The
empirical background uses 1,916 WikiText word tokens. These baselines are not
matched to the 1,000-token sample. A paired test of BM25 against the hybrid
is not significant on either collection.}
\label{tab:retrieval-controls}
\end{table}

The LM prior helps only against a very sparse empirical background. At 1,000
reference tokens, the hybrid beats empirical counts by .051 nDCG@10 on SciFact
(.650 versus .599, 95\% paired query CI [.026, .077], adjusted $p=.0068$
across the 68-test family) and by .018 on NFCorpus (CI [.008, .029], adjusted
$p=.0312$). Neither pure LM comparison survives correction. At 10,000 tokens
and above, the hybrid's advantage disappears or reverses. The full-data
empirical prior reaches .656 on SciFact, above the 1,000-token hybrid.

The standard baselines are as good or better. On SciFact, BM25 reaches .663,
collection Dirichlet .659, and Jelinek-Mercer .657, all above the hybrid's
.650. Even a background of 1,916 WikiText tokens mixed with collection
frequencies reaches .661 and beats the LM hybrid. None of these differences is
reliable. BM25 minus the hybrid is +.013 nDCG@10 (95\% CI [$-.005$, .032],
$p=.146$, Holm-adjusted $p=.292$). The collection Dirichlet, Jelinek-Mercer,
and WikiText-mix differences from the hybrid (+.009, +.007, and +.011) have
adjusted $p=1.0$ in the 40-contrast control family. On NFCorpus, the hybrid's
.308 is close to BM25 at .305 (BM25 minus hybrid $-.003$, CI [$-.011$, .005],
adjusted $p=.465$) and above collection Dirichlet at .298. The collection Dirichlet difference ($-.010$,
CI [$-.019$, $-.001$]) does not survive correction (adjusted $p=.45$).
Jelinek-Mercer is lower than the hybrid by .016 (adjusted $p=.036$). Uniform
and additive smoothing are worse than all other methods on both collections.
The LM prior is therefore a fix for an artificially starved background, not
an improvement over standard practice.

The direct scarcity test compares each prior's gain at 1,000 tokens with its
gain at the largest feasible sampled budget. On SciFact, the LM and hybrid
gains decline by .054 and .067 nDCG (95\% CIs [.030, .078] and [.043, .092],
adjusted $p=.0068$ for each). On NFCorpus, the hybrid gain declines by .018
(CI [.008, .027], adjusted $p=.0165$), while the LM-only decline of .008 is
uncertain (CI [$-.001$, .017], adjusted $p=1.0$).

Held-out token cross-entropy ranks the priors differently from retrieval. At
1,000 tokens, LM cross-entropy is 11.82 on SciFact and 12.09 on NFCorpus,
compared with empirical values of 7.30 and 7.59. The hybrid is intermediate.
Lower query cross-entropy therefore does not explain the hybrid's retrieval
gain.

Two further studies test other uses of LM statistics in BM25. LM-derived IDF
did not improve nDCG@10 over BM25 on SciFact, NFCorpus, or FiQA. Hybrid IDF was
statistically indistinguishable from BM25 on SciFact and NFCorpus and worse on
FiQA. Adaptive LM weighting reduced nDCG@10 on all three collections. Each LM
method beat its formula-matched empirical-prior control, but not BM25. A prior
can win inside a chosen formula and still lose to the standard ranker. In the
stop-word study, rank-frequency breakpoints were fitted to collection document
frequency, token frequency, frozen model probabilities, and external empirical
frequencies. The maximum of 705 terms was selected for the first three signals
in every collection, while external empirical frequency selected 305. LM-based
filtering improved nDCG@10 over document-frequency filtering on all three
collections, but lost to external empirical filtering and standard BM25 on
SciFact and FiQA, with no reliable advantage over BM25 on NFCorpus. Some
filtered queries returned no documents, up to 59 of 323 NFCorpus queries in
one source-target condition.

\subsection{RQ1: Stylometry and lexical tasks}

\begin{table}[t]
\centering
\small
\begin{tabular}{llrrr}
\toprule
Corpus & Tokens & Empirical & LM & Hybrid \\
\midrule
Blogs & 1,000 & .715 & .716 & .711 \\
Blogs & 10,000 & .713 & .716 & .711 \\
Blogs & Full & .712 & .716 & .712 \\
Novels & 1,000 & .619 & .625 & .625 \\
Novels & 10,000 & .620 & .625 & .624 \\
Novels & Full & .620 & .625 & .621 \\
\bottomrule
\end{tabular}
\caption{Global pairwise AUC for stylometric attribution. Sampled results
average three seeds. No LM or hybrid comparison survives Holm correction
over the 18 stylometry contrasts.}
\label{tab:stylometry}
\end{table}

The planned stylometry family has 18 contrasts: LM minus empirical, hybrid
minus empirical, and hybrid minus LM, for two corpora at 1,000 tokens, 10,000
tokens, and full data. Holm correction covers these 18 tests. Neither the LM
prior nor the hybrid reliably improves over empirical counts. Full-data AUC gains for the LM
range from .004 to .005, and cluster-bootstrap intervals include zero. Two
hybrid-minus-empirical differences are nominally significant before
correction: blogs at 10,000 tokens ($-.001$, $p=.022$, 18-test adjusted
$p=.40$) and novels at full data (+.001, $p=.037$, 18-test adjusted $p=.64$). Neither survives
correction, and they point in opposite directions. The 1,000-to-10,000-token stylometry changes are small, and all
four author-cluster intervals include zero, so the retrieval scarcity pattern
does not carry over. Top-one attribution is mixed. At full data, the LM scores
.280 versus .293 for empirical counts on blogs, and .175 versus .150 on novels.
The hybrid matches the empirical top-one scores on both full-data conditions.

Held-out cross-entropy does not track attribution either. At full reference
size, cross-entropy for empirical, LM, and hybrid priors is 4.651, 4.933, and
4.651 on blogs, and 4.678, 4.847, and 4.695 on novels. At 1,000 tokens, the LM
stays at 4.933 and 4.847, empirical counts give 4.814 and 4.838, and the hybrid
gives 4.842 and 4.833. The LM prior alone does not improve held-out word
prediction on these collections.

\begin{table*}[t]
\centering
\small
\begin{tabular}{lrrrrrrrrr}
\toprule
Prior & \multicolumn{3}{c}{Distinctiveness} &
\multicolumn{3}{c}{Next-word prediction} &
\multicolumn{3}{c}{Spelling correction} \\
 & 1,000 & 10,000 & 100,000 & 1,000 & 10,000 & 100,000 &
  1,000 & 10,000 & 100,000 \\
\midrule
Empirical & .526 & .535 & .505 & .269 & .359 & .465 & .855 & .897 & .916 \\
LM & .518 & .503 & .510 & .401 & .401 & .401 & .902 & .902 & .902 \\
Hybrid & .556 & .545 & .514 & .401 & .403 & .408 & .902 & .909 & .918 \\
Smoothed bigram & n/a & n/a & n/a & .336 & .418 & .505 & n/a & n/a & n/a \\
Full frequency & .522 & .526 & .516 & .636 & .636 & .636 & .925 & .925 & .925 \\
Uniform & .543 & .509 & .497 & .215 & .215 & .215 & .785 & .785 & .785 \\
\bottomrule
\end{tabular}
\caption{Top-one accuracy pooled over items and three reference-sampling
seeds. Accuracy levels in the text use this pooled estimator. Paired
differences and their intervals use author-cluster means, so they can differ
slightly from differences of table values.
Distinctiveness classifies held-out blog authors from selected terms.
Next-word prediction and spelling use constructed candidate sets. Next-word
prediction is the LM training objective. Full-frequency values use all
training text.}
\label{tab:extensions}
\end{table*}

In lexical distinctiveness, the top-100-term classifier does not reliably
improve over the empirical-prior classifier (Table~\ref{tab:extensions}). At
1,000 tokens, the LM accuracy difference is $-.012$ (95\% author-cluster CI
[$-.051$, .030]) and the hybrid difference is .020 (CI [$-.029$, .072]).
Neither survives correction. The hybrid lowers token cross-entropy from 8.03 to
7.66 at this budget, so better word prediction again does not give a reliable
task gain.

Next-word prediction is the language model's training objective, so a gain
there shows cloze ability and nothing more. At 1,000 tokens, the LM and hybrid
improve top-one accuracy over sampled pair counts by .131 (95\% author-cluster
CI [.102, .159], adjusted $p=.0068$). The advantage shrinks with more counts.
At 100,000 tokens, sampled pair counts reach .465, above the LM at .401 and the
hybrid at .408, and full-reference frequency reaches .636 at all budgets. In
the post-hoc comparison with the smoothed bigram, the LM and hybrid lead by
.065 at 1,000 tokens (CIs [.039, .091], adjusted $p=.0165$), are not reliably
different at 10,000 tokens, and trail by .104 and .096 at 100,000 tokens (CIs
[$-.129$, $-.079$] and [$-.120$, $-.073$], adjusted $p=.0068$). PMI and the
t-score reach .192 and .239 at 1,000 tokens and .337 and .436 at 100,000.
Empirical candidate cross-entropy is lower than LM cross-entropy at both 1,000
tokens (1.55 versus 1.70) and 100,000 tokens (1.23 versus 1.70). This task says
nothing about human collocation judgments.

In synthetic spelling correction, LM and hybrid priors improve top-one
accuracy over sampled unigram counts at 1,000 tokens by .045 and .046
(author-cluster means, 95\% CIs [.036, .054] and [.039, .052], adjusted $p=.0068$ for both).
The gains shrink as counts grow, by .063 for the LM and .044 for the hybrid
(both adjusted $p=.0068$). At 100,000 tokens, the hybrid is .918 versus .916
for empirical counts, the LM is .902, and full-reference frequency is .925.
Typos are generated from the known answer with the same edit channel that
supplies candidates, so this is a sanity check, not evidence of utility for
natural spelling errors.

Estimating the retrieval vocabulary's LM prior took 333.6 seconds on Apple
MPS, once for both collections and all budgets. Downstream scoring, tuning,
and reporting took 93.1 seconds on CPU. Extension word and pair scores took
537.7 seconds on an NVIDIA GB10, 591.6 seconds on Apple MPS, and 744.3 seconds
on an AMD Radeon 8060S, as separate shard runtimes. The ROCm shard overlapped
with another GPU workload. I did not measure query-time latency.

\subsection{RQ2: Close-language identification}

\begin{table*}[t]
\centering
\small
\begin{tabular}{lrrrrrr}
\toprule
 & \multicolumn{3}{c}{Raw (planned)} & \multicolumn{3}{c}{Dev-calibrated (post hoc)} \\
System & All four & bos/hrv & ind/zsm & All four & bos/hrv & ind/zsm \\
\midrule
Causal model, from scratch & 77.94 & 66.21 & 89.67 & 81.74 & 68.63 & 94.91 \\
Unigram, same subword tokenizer & 76.24 & 59.88 & 92.59 & 79.72 & 66.11 & 93.33 \\
UniLID (\texttt{unilid-1940}) & n/a & n/a & 86.31 & n/a & n/a & n/a \\
fastText lid.176 & 62.99 & 55.53 & 70.55 & 71.20 & 61.81 & 80.68 \\
GlotLID-M & 82.19 & 70.75 & 93.63 & 86.34 & 76.38 & 96.29 \\
\midrule
Character 1-5 gram naive Bayes (post hoc) & 73.07 & 58.55 & 87.60 & 79.27 & 66.35 & 92.24 \\
\bottomrule
\end{tabular}
\caption{FLORES-200 devtest accuracy (\%) on 4,048 full sentences. ``All four''
is four-way classification. Pair columns restrict the decision to the two
languages of the pair, for sentences of that pair. The raw evaluation was the
planned final evaluation. The dev-offset calibration was added afterward, and
each system's offsets are fitted on the same 3,988 dev sentences. The
character model was added later still. The released UniLID model has no
Bosnian or Croatian label and was run uncalibrated on Indonesian/Malay only.
The causal, unigram, and character models use the same 50 million training
tokens per language. One training seed.}
\label{tab:lid}
\end{table*}

A separate study trained one causal Transformer per language from scratch on
Bosnian, Croatian, Indonesian, and Malay Wikipedia \citep{wikimedia2023} and
scored FLORES-200 devtest sentences \citep{nllb2022} (Table~\ref{tab:lid}).
The matched unigram shares the causal model's 16,000-type subword tokenizer
and training tokens. The planned comparators were this unigram, fastText
lid.176, GlotLID-M, and the released UniLID model
\citep{meister2026unilid}. UniLID's label set covers Indonesian and Malay but
not Bosnian or Croatian, so it was scored on that pair only, without
calibration.

The planned raw evaluation is mixed. The causal model reaches 77.94\% overall,
above the unigram at 76.24\% and below GlotLID-M at 82.19\%. It beats the
unigram on Bosnian/Croatian by 6.32 points but loses on Indonesian/Malay,
89.67\% versus 92.59\% ($-2.92$ points, 95\% paired CI [$-4.40$, $-1.43$]).
On Indonesian/Malay, UniLID reaches 86.31\%, below the causal model, the
matched unigram, and GlotLID-M.

Two analyses are post hoc. First, after the raw evaluation I fitted additive
per-language score offsets for each system on FLORES dev and froze them for
devtest. Second, because the matched unigram can carry the shared tokenizer's
biases, I added a character 1-5 gram naive Bayes model that uses no subword
tokenizer. It counts hashed character n-grams ($2^{22}$ buckets per order) in
the same training tokens, decoded back to text. Its additive smoothing was
selected on dev from 0.01, 0.1, and 1.0, and its offsets were fitted on the
same dev split. All intervals below are 95\% paired bootstrap intervals over
FLORES sentences, without multiplicity correction across the nine calibrated
contrasts. The bos/hrv lower bounds of 0.20 and 0.40 points are the most
fragile.

After the post-hoc calibration, the causal model leads the unigram by 2.03
points overall (CI [0.77, 3.26]) and by 2.52 and 1.58 points on the two pairs.
It leads the character n-gram model by 2.47 points overall (CI [1.24, 3.71]),
2.27 on Bosnian/Croatian (CI [0.20, 4.35]), and 2.67 on Indonesian/Malay (CI
[1.48, 3.90]). The character model does not differ reliably from the subword
unigram after calibration ($-0.44$ points, CI [$-1.51$, 0.57]). Raw, the
character model is the weaker of the two (73.07\% versus 76.24\%). Under the
post-hoc calibration, the contextual gain is not an artifact of the shared
tokenizer, at least against this simple character baseline.

The released specialist is still better. GlotLID-M leads the causal model by
4.59 points after calibration (CI [3.36, 5.88]) and by 4.25 points raw, and by
7.76 points on Bosnian/Croatian after calibration. fastText lid.176 trails all
other systems. The causal gain is therefore a bounded measurement of context
value in two close pairs. It does not show that contextual likelihood beats
strong identifiers, and the models are trained from scratch, so it says
nothing about pretrained base models. The character baseline is a simple
hashed naive Bayes model, not a strong system. Its smoothing was selected at
the edge of its grid, although dev accuracy changed by only 0.33 points across
the 100-fold range. A tuned character model, such as the SVM and
logistic-regression systems of the DSL shared tasks, or OpenLID-v3 may score
higher.

\subsection{RQ3: Existing human-judgment data}

I piloted five tasks. Four came from a screen of 31 candidate tasks, and I
added EIFFEL beside that screen as a related idiom-completion task. The screen
dropped candidates without a clear data license, a strong classical baseline,
or a context-dependent target. The five tasks tested whether probabilities
from Qwen3.5-2B-Base and Qwen3.5-4B-Base \citep{qwen35} add predictive value
beyond classical features.
Each analysis is five-fold cross-validation, grouped by verb, expression,
prompt, or target. DAIS, EIFFEL, and MERLIN use all labeled items. SWORDS uses
100 of its 762 test targets, and the idiomaticity pilot uses a 2,000-row
sample. The same folds were used to select the pilots, so no independent
confirmatory split remains. A Holm correction over the five candidates keeps
the three positive results (adjusted $p\le.001$, the floor of 10,000
bootstrap resamples). For DAIS, this gate is
the post-hoc additive gate described below. With 31 screened candidates, a
correction over only the five piloted tasks understates the selection.

The DAIS dative-preference task \citep{hawkins2020dais} first failed its
registered test. Hawkins et al.\ introduced DAIS and showed that neural
language-model probabilities predict human dative preferences. My registered
pilot asked whether boundary-corrected whole-word scores alone track the
human preference better than the strongest classical model. They did not. The
held-out-verb Spearman correlation was .527 for the 2B model and .474 for the
4B model, against .627 for the feature baseline (differences $-.100$ and
$-.153$, verb-cluster CIs [$-.142$, $-.058$] and [$-.198$, $-.108$]). After
this failure, I changed the success criterion to an additive test. That change
is post hoc. In the additive test, the same saved scores reduced
verb-held-out MSE beyond a 5-gram and feature model by 8.87\% (2B) and 8.44\%
(4B). The target is a continuous preference from 0 to 100. The additive
result is consistent with Hawkins et al.\ and is a replication under a
post-hoc criterion, not a new human-judgment finding.

EIFFEL \citep{eiffel} asks for the intended completion of a French idiom.
Adding LM scores reduced log loss by 9.46\% (2B) and 20.81\% (4B) beyond
character and word n-gram features. Idiom completion is a cloze task, which is
the language-modeling objective, and surprisal already explains part of the
human idiom processing advantage \citep{socolof2025idiom}. The gain shows
cloze ability, not linguistic utility beyond that objective. MERLIN
\citep{merlin2014} German learner proficiency is less directly aligned with
the objective. Mean corrected surprisal reduced CEFR prediction MAE by
16.32\% and 17.16\% beyond word and character n-grams. Its cluster bootstrap
resamples only 16 exam-prompt clusters, so its intervals may understate
variation across prompts. This result is not new. GPT perplexity already
tracks CEFR level in Swedish learner essays \citep{munozsanchez2024perplexity},
and perplexity features from LLMs already serve as CEFR classifiers across
learner corpora \citep{stearns2026cefr}. My pilot extends that idea to German
MERLIN texts with an additive test against n-gram features. It is not a
training-objective result and not a Hawkins et al.\ replication.

The full-context SWORDS lexical-substitution pilot \citep{lee2021swords} and
the French Lexical Network idiomaticity-class pilot \citep{liu2024idiom} did
not pass. SWORDS was first scored from left context only, while
the human judgments and the classical baseline used the full context. I
rescored the same 100 of 762 target contexts with the full context, and that
rescored version is the one reported. It has too little power for the planned
effect on the full test. The first idiomaticity scoring run stopped at a sentence-initial target,
which a causal model cannot score. I restricted the sample to targets with
preceding context and rescored it. The pilot used a balanced
2,000-item sample and showed no reliable additive value. Its labels are the
four expert lexicon classes of the French Lexical Network, not human
judgments. The three positive
and two negative pilots are separate task-specific results, not a general
claim that pretrained probabilities improve prediction of human judgments.

\subsection{Task-trained models and lexical laws}

A matched-pretraining study compared standard training (A), linguistically
curated training (B), and curated training with a pair-ranking objective (C)
on one compact GPT-2 model with one training seed. All three used 24,462
updates from a shared initialization. I do not report its effects. The file
of paired effects was lost, and the saved checkpoints have not been rescored.
This paper makes no claim from this study.

A descriptive study compared frozen whole-word probabilities from DistilGPT-2
and GPT-2 \citep{radford2019gpt2} with human frequencies, and measured Zipf
fits, Heaps vocabulary growth, and Taylor burstiness in held-out human text and
generated samples from WikiText-2 and Project Gutenberg domains. A
100-replicate bootstrap rejected the pure Zipf model at $p=.0099$, the
smallest value 100 replicates can give, for all 14
human and generated count distributions. Human held-out Heaps exponents were
.734 for Wikipedia and .664 for Gutenberg. At a fixed 20,000-word generation
budget, generated samples contained 285 to 13,201 types. Temperature 1.0 had
the lowest Jensen-Shannon divergence from human text in all four model-domain
pairs. Human Taylor exponents ranged from 1.00 to 1.11, and generated
exponents from .97 to 1.51. Earlier work already measured these laws in
neural and LLM text, including temperature effects
\citep{takahashi2017neural,takahashi2019scaling,mikhaylovskiy2025zipf}. These
numbers describe two English domains and two GPT-2-family models. They are not
new laws, and they do not show that either distribution is a better baseline
for any task.

\section{Analysis}

The evidence supports a narrow conclusion. A frozen LM word distribution
repairs an artificially starved retrieval background, but it does not beat BM25 or
collection Dirichlet smoothing on SciFact, and it is only close to BM25 on
NFCorpus. It gives no reliable gain in stylometry or distinctive-term
classification, and LM-derived IDF and stop words do not beat BM25. In
language identification, context helps beyond a matched unigram and a simple
character model only after post-hoc calibration, and a released identifier is
better in both the raw and the calibrated analysis.

The clearest positive results come from tasks that match the language-modeling
objective: next-word prediction and idiom completion. These show that a language model predicts held-out words
well. They do not show that its probabilities measure anything beyond that.
The dative result repeats a known finding. MERLIN is the one positive pilot
that is less tied to the objective. It has only 16 prompt clusters, and prior
work already links perplexity to CEFR level.

Predictive fit and instrument utility also diverge. Lower held-out
cross-entropy did not explain the retrieval gain, did not track stylometric
AUC, and did not give a reliable distinctiveness gain. Predictability, rarity,
retrieval relevance, and burstiness are related but different quantities. A
model can estimate one well without being a good proxy for the others.

\section{Conclusion}

Across these studies, language-model distributions sometimes improve a defined
comparison, but the gains depend on task, data budget, and training objective.
In retrieval, the gain holds only against a background from 1,000 sampled
tokens, and the standard baselines remain as good or better. In close-language
identification, a contextual model beats a matched unigram and a character
n-gram model after post-hoc calibration, but GlotLID-M is better. The positive human-judgment results
align with the cloze objective or repeat prior work. Each language-model distribution
is a candidate measurement whose value must be shown against a task-matched
baseline. The evidence does not support it as a general replacement for corpus
statistics.

\section{Limitations}

The studies do not form one matched experiment. The pretrained-prior
evaluations mostly use one frozen DistilGPT-2 checkpoint and 32 WikiText-103
contexts, the lexical-law study uses two GPT-2-family checkpoints, and the
language-identification study
has one training seed and four close languages. The human-judgment pilots use
different datasets and cross-validation folds that also selected the pilots,
with no independent confirmatory split. These differences prevent a pooled
estimate and limit transfer across model families and tasks.

The retrieval evidence covers two collections for the prior comparison and
three for the BM25 studies. It does not compare with neural rankers such as
SPLADE or LLM query-likelihood models. The sampled-budget setting is artificial
for query likelihood, which already reads the full collection. The hybrid and
Dirichlet parameters are chosen on small development sets, while BM25 keeps
its default $k_1=1.2$ and $b=0.75$. A tuned BM25 could score higher.

The character n-gram baseline is one simple hashed naive Bayes model with
three smoothing values, and the selected value lies at the grid edge. It is a
control for the tokenizer, not the strongest character model. The
language-identification calibration and the character control were both added
after the raw evaluation, so they are post hoc. UniLID covers only one of the
two pairs, and OpenLID was not run. The blog job
labels are coarse. Task and domain vary together across datasets, so the
cross-task differences do not identify domain-shift effects. The next-word
task has no human collocation labels. Spelling uses synthetic single-edit
errors and a uniform channel. The two novel tasks have 20 author clusters.
Tests account for query and author clustering and planned comparisons, but not
for unmeasured models, datasets, or evaluation choices.

\bibliography{references}
\end{document}
